\documentclass[12pt,a4paper]{report}

% ── Packages ──────────────────────────────────────────────────────────────────
\usepackage[utf8]{inputenc}
\usepackage[T1]{fontenc}
\usepackage{lmodern}
\usepackage[margin=2.5cm]{geometry}
\usepackage{amsmath,amssymb,amsthm}
\usepackage{mathtools}
\usepackage{bm}           % bold math
\usepackage{graphicx}
\usepackage{booktabs}
\usepackage{array}
\usepackage{longtable}
\usepackage{multirow}
\usepackage{xcolor}
\usepackage{hyperref}
\usepackage{listings}
\usepackage{caption}
\usepackage{subcaption}
\usepackage{float}
\usepackage{enumitem}
\usepackage{titlesec}
\usepackage{fancyhdr}
\usepackage{tikz}
\usetikzlibrary{shapes,arrows,positioning,fit,backgrounds}

% ── Hyperref setup ────────────────────────────────────────────────────────────
\hypersetup{
    colorlinks=true,
    linkcolor=blue!60!black,
    urlcolor=blue!60!black,
    citecolor=green!50!black,
}

% ── Header / footer ───────────────────────────────────────────────────────────
\pagestyle{fancy}
\fancyhf{}
\fancyhead[L]{\nouppercase{\leftmark}}
\fancyhead[R]{\thepage}
\renewcommand{\headrulewidth}{0.4pt}

% ── Code listing style ────────────────────────────────────────────────────────
\lstset{
  basicstyle=\ttfamily\small,
  keywordstyle=\color{blue!70!black}\bfseries,
  commentstyle=\color{green!50!black}\itshape,
  stringstyle=\color{red!60!black},
  numberstyle=\tiny\color{gray},
  numbers=left,
  frame=single,
  breaklines=true,
  showstringspaces=false,
  tabsize=4,
}

% ── Math helpers ──────────────────────────────────────────────────────────────
\newcommand{\vect}[1]{\bm{#1}}
\newcommand{\mat}[1]{\mathbf{#1}}
\newcommand{\R}{\mathbb{R}}
\newcommand{\skew}[1]{[#1]_{\times}}
\DeclareMathOperator{\diag}{diag}
\DeclareMathOperator*{\argmin}{arg\,min}

% ── Title ─────────────────────────────────────────────────────────────────────
\title{
    \vspace{-1cm}
    {\Large\textbf{Dynamic Locomotion through Convex Optimisation}}\\[0.8em]
    {\large Complete Technical Documentation}\\[0.4em]
    {\normalsize Convex MPC Trot Controller for the Unitree Go2 Quadruped Robot}\\[0.3em]
    {\normalsize Flat-Ground Trotting $\cdot$ Stair Ascent \& Descent $\cdot$ Terrain Adaptation}\\[1em]
    {\small Based on: Di Carlo et al., ``Dynamic Locomotion in the MIT Cheetah 3\\
    Through Convex Model-Predictive Control'', IROS 2018}
}
\author{Stanny}
\date{\today}

% ══════════════════════════════════════════════════════════════════════════════
\begin{document}
% ══════════════════════════════════════════════════════════════════════════════

\maketitle
\tableofcontents
\listoftables
\newpage

% ══════════════════════════════════════════════════════════════════════════════
\chapter{Introduction and Project Overview}
% ══════════════════════════════════════════════════════════════════════════════

\section{Motivation}

Legged robots must balance three competing requirements simultaneously: they must
generate sufficient ground-reaction forces (GRFs) to support their weight,
satisfy friction constraints at each foot, and coordinate leg swing trajectories
that advance the body forward.  Classical approaches couple these problems,
making real-time control difficult.

This project implements the approach of Di~Carlo et al.\ \cite{diCarlo2018},
which decouples stance-force optimisation from swing-foot placement by
formulating the stance problem as a \emph{convex} Quadratic Programme (QP)
that can be solved on a $30\,\mathrm{ms}$ horizon in real time.

\section{Platform}

\textbf{Hardware:} Unitree Go2 quadruped robot.\\
\textbf{Simulation:} MuJoCo 3.x (via the \texttt{unitree\_mujoco} scene).\\
\textbf{Solver:} \texttt{quadprog} (primary) / \texttt{osqp} (fallback).

\section{Feature Summary}

\begin{itemize}[noitemsep]
    \item Flat-ground diagonal trot at commanded forward velocity.
    \item Stair ascent and descent over a 5-step staircase
          ($6\,\mathrm{cm}$ riser, $28\,\mathrm{cm}$ tread, $30\,\mathrm{cm}$ total height).
    \item Adaptive CoM height, body pitch reference, and swing clearance.
    \item Terrain-aware friction-cone constraints via support-plane estimation.
    \item Per-foot stair tread snapping and early-touchdown detection.
    \item Collapse detection and PD recovery mode.
\end{itemize}

\section{Repository Structure}

\begin{center}
\begin{tabular}{ll}
\toprule
\textbf{File} & \textbf{Role} \\
\midrule
\texttt{robot\_params.py}     & Physical parameters, Q weights, RobotState dataclass \\
\texttt{config.py}            & MuJoCo IDs, hardware limits, debug flags \\
\texttt{dynamics.py}          & SRB A/B matrices, ZOH discretisation \\
\texttt{gait\_scheduler.py}   & Trot gait contact schedule \\
\texttt{state\_estimator.py}  & Ground-truth reader, contact detection \\
\texttt{mpc\_solver.py}       & Condensed QP: condense, cost, constraints, solve \\
\texttt{swing\_controller.py} & Minimum-jerk swing trajectory, task-space PD \\
\texttt{torque\_utils.py}     & GRF$\to$torques, reference builder, PD stand, logger \\
\texttt{terrain\_estimator.py}& Stair geometry, IIR terrain tracking \\
\texttt{support\_plane.py}    & SVD plane fit of stance feet \\
\texttt{mpc\_controller.py}   & Phase A/B/C state machine (top-level) \\
\texttt{simulation.py}        & MuJoCo viewer + headless runner \\
\texttt{flat\_trot.py}        & Entry point — flat ground \\
\texttt{stairs\_climb.py}     & Entry point — stair scene \\
\bottomrule
\end{tabular}
\end{center}

% ══════════════════════════════════════════════════════════════════════════════
\chapter{Robot Model and Physical Parameters}
% ══════════════════════════════════════════════════════════════════════════════

\section{Go2 Physical Parameters}

The Unitree Go2 is modelled as a \emph{Single Rigid Body} (SRB) at the
centroidal level.  Only the trunk mass and inertia enter the MPC; leg masses
(approx.\ $1.5\,\mathrm{kg}$/leg) are compensated by adding the joint-space
gravity/Coriolis bias $\vect{q}_{\text{frc\_bias}}$ after the QP solve.

\begin{table}[H]
\centering
\caption{Go2 physical parameters used in the MPC}
\begin{tabular}{lll}
\toprule
\textbf{Parameter} & \textbf{Symbol} & \textbf{Value} \\
\midrule
Total mass            & $m$      & $15.206\,\mathrm{kg}$ \\
Gravity               & $g$      & $9.81\,\mathrm{m/s^2}$ \\
Roll inertia (body)   & $I_{xx}$ & $0.10703\,\mathrm{kg\,m^2}$ \\
Pitch inertia (body)  & $I_{yy}$ & $0.09808\,\mathrm{kg\,m^2}$ \\
Yaw inertia (body)    & $I_{zz}$ & $0.02445\,\mathrm{kg\,m^2}$ \\
Hip ab/ad offset      & $l_\mathrm{hip}$   & $0.08085\,\mathrm{m}$ \\
Thigh link length     & $l_\mathrm{thigh}$ & $0.213\,\mathrm{m}$ \\
Calf link length      & $l_\mathrm{calf}$  & $0.213\,\mathrm{m}$ \\
\midrule
Hip origin FL (body frame) & & $(+0.1934,\,+0.1125,\,0)$ \\
Hip origin FR (body frame) & & $(+0.1934,\,-0.1125,\,0)$ \\
Hip origin RL (body frame) & & $(-0.1934,\,+0.1125,\,0)$ \\
Hip origin RR (body frame) & & $(-0.1934,\,-0.1125,\,0)$ \\
\bottomrule
\end{tabular}
\end{table}

\section{Body-Frame Inertia Tensor}

\begin{equation}
    \mat{I}_B = \diag\!\left(I_{xx},\; I_{yy},\; I_{zz}\right)
              = \diag(0.1070,\; 0.0981,\; 0.0245)\;\mathrm{kg\,m^2}
\end{equation}

\section{MPC State Vector}

The 13-dimensional state vector used throughout the MPC is:

\begin{equation}
    \vect{x} = \begin{bmatrix}
        \phi & \theta & \psi & p_x & p_y & p_z &
        \omega_x & \omega_y & \omega_z & v_x & v_y & v_z & {-g}
    \end{bmatrix}^\top \in \R^{13}
    \label{eq:state}
\end{equation}

\begin{table}[H]
\centering
\caption{MPC state vector components}
\begin{tabular}{cllll}
\toprule
\textbf{Index} & \textbf{Symbol} & \textbf{Meaning} & \textbf{Unit} & \textbf{Frame} \\
\midrule
0  & $\phi$     & Roll  (Euler ZYX)      & rad    & World \\
1  & $\theta$   & Pitch (Euler ZYX)      & rad    & World \\
2  & $\psi$     & Yaw   (Euler ZYX)      & rad    & World \\
3  & $p_x$      & CoM position $x$       & m      & World \\
4  & $p_y$      & CoM position $y$       & m      & World \\
5  & $p_z$      & CoM position $z$       & m      & World \\
6  & $\omega_x$ & Angular velocity $x$   & rad/s  & World \\
7  & $\omega_y$ & Angular velocity $y$   & rad/s  & World \\
8  & $\omega_z$ & Angular velocity $z$   & rad/s  & World \\
9  & $v_x$      & Linear velocity $x$    & m/s    & World \\
10 & $v_y$      & Linear velocity $y$    & m/s    & World \\
11 & $v_z$      & Linear velocity $z$    & m/s    & World \\
12 & $-g$       & Gravity augmentation   & m/s$^2$ & — \\
\bottomrule
\end{tabular}
\end{table}

The \emph{gravity augmentation} trick appends the constant $-g$ as the 13th
state so that gravity appears naturally in the linear dynamics via the
$\mat{A}$ matrix without needing an affine term.

% ══════════════════════════════════════════════════════════════════════════════
\chapter{Single Rigid Body Dynamics}
% ══════════════════════════════════════════════════════════════════════════════

\section{Centroidal Model Motivation}

A quadruped with $4 \times 3 = 12$ leg joints has 18 degrees of freedom
(6 body + 12 joints).  Real-time whole-body optimisation over a predictive
horizon is intractable.  The \emph{single rigid body} (SRB) model reduces the
problem to 13 states and $4\times 3 = 12$ inputs (one 3-D force vector per
foot), giving a convex QP solvable in $<5\,\mathrm{ms}$.

The SRB assumptions are:
\begin{enumerate}[noitemsep]
  \item Leg masses are small compared to the trunk ($\approx 10\%$ each).
  \item Foot accelerations are small (slow enough swings).
  \item The inertia tensor can be approximated by rotating the body-frame
        tensor by yaw only (Di~Carlo approximation).
\end{enumerate}

\section{Newton--Euler Equations}

For a rigid body with $n_c$ foot contacts, Newton's second law gives:

\begin{align}
    m\,\dot{\vect{v}} &= \sum_{i=1}^{n_c} \vect{f}_i + m\,\vect{g}
    \label{eq:newton} \\[4pt]
    \hat{\mat{I}}\,\dot{\vect{\omega}} &=
        \sum_{i=1}^{n_c} \vect{r}_i \times \vect{f}_i
    \label{eq:euler}
\end{align}

where $\vect{f}_i \in \R^3$ is the ground-reaction force (GRF) at foot $i$,
$\vect{r}_i = \vect{p}_i - \vect{p}_\mathrm{CoM}$ is the moment arm from the
CoM to foot $i$, and $\hat{\mat{I}}$ is the world-frame inertia tensor.

\subsection{Di Carlo Inertia Approximation}

Computing the full world-frame inertia $\mat{R}\,\mat{I}_B\,\mat{R}^\top$
at every time step introduces trilinear products that destroy convexity.
Di~Carlo et al.\ approximate this by rotating only around yaw:

\begin{equation}
    \hat{\mat{I}}(\psi) = \mat{R}_z(\psi)\;\mat{I}_B\;\mat{R}_z(\psi)^\top
    \label{eq:inertia_approx}
\end{equation}

where $\mat{R}_z(\psi)$ is the rotation matrix for yaw angle $\psi$.
This is exact on flat ground when roll and pitch are small, which is a
reasonable approximation during trotting locomotion.

\section{Continuous-Time State-Space Model}

\subsection{Euler-Angle Kinematics}

Using the small-angle Euler-rate approximation from Eq.~(12) of Di~Carlo:

\begin{equation}
    \begin{bmatrix}\dot\phi \\ \dot\theta \\ \dot\psi\end{bmatrix}
    \approx \mat{R}_z(\psi)^\top\;\vect{\omega}
\end{equation}

This gives the upper-left block of $\mat{A}_c$:
\begin{equation}
    \mat{A}_c[0\!:\!3,\; 6\!:\!9] = \mat{R}_z(\psi)^\top
\end{equation}

\subsection{Position Kinematics}

Linear velocity integrates directly into position:
\begin{equation}
    \mat{A}_c[3\!:\!6,\; 9\!:\!12] = \mat{I}_3
\end{equation}

\subsection{Gravity Augmentation}

The 13th state $x_{12} = -g$ feeds into $\dot v_z$:
\begin{equation}
    \mat{A}_c[11,\; 12] = 1
    \quad\Longrightarrow\quad
    \dot v_z \supseteq x_{12} = -g
\end{equation}

\subsection{Full Continuous A Matrix}

\begin{equation}
    \mat{A}_c(\psi) =
    \begin{pmatrix}
        \mat{0}_{3\times3} & \mat{0}_{3\times3} & \mat{R}_z(\psi)^\top & \mat{0}_{3\times3} & \mat{0}_{3\times1} \\
        \mat{0}_{3\times3} & \mat{0}_{3\times3} & \mat{0}_{3\times3} & \mat{I}_3 & \mat{0}_{3\times1} \\
        \mat{0}_{3\times3} & \mat{0}_{3\times3} & \mat{0}_{3\times3} & \mat{0}_{3\times3} & \mat{0}_{3\times1} \\
        \mat{0}_{3\times3} & \mat{0}_{3\times3} & \mat{0}_{3\times3} & \mat{0}_{3\times3} & \vect{e}_3 \\
        \mat{0}_{1\times3} & \mat{0}_{1\times3} & \mat{0}_{1\times3} & \mat{0}_{1\times3} & 0
    \end{pmatrix}
    \in \R^{13\times13}
    \label{eq:Ac}
\end{equation}

where $\vect{e}_3 = [0,0,1]^\top$ places the gravity constant in the $v_z$ equation.

\subsection{Continuous B Matrix}

The $\mat{B}_c$ matrix maps the control input
$\vect{u} = [\vect{f}_1^\top, \vect{f}_2^\top, \vect{f}_3^\top, \vect{f}_4^\top]^\top \in \R^{12}$
to state derivatives.  For each stance foot $i$ ($c_i = 1$):

\begin{equation}
    \mat{B}_c[6\!:\!9,\; 3i\!:\!3(i+1)] = \hat{\mat{I}}^{-1}\,\skew{\vect{r}_i}
    \qquad\text{(angular dynamics)}
    \label{eq:Bc_ang}
\end{equation}

\begin{equation}
    \mat{B}_c[9\!:\!12,\; 3i\!:\!3(i+1)] = \frac{\mat{I}_3}{m}
    \qquad\text{(linear dynamics)}
    \label{eq:Bc_lin}
\end{equation}

where $\skew{\vect{r}}$ denotes the $3\times3$ skew-symmetric matrix
(cross-product operator) of vector $\vect{r}$:

\begin{equation}
    \skew{\vect{r}} =
    \begin{pmatrix}
        0 & -r_z & r_y \\
        r_z & 0 & -r_x \\
        -r_y & r_x & 0
    \end{pmatrix}
\end{equation}

Swing legs ($c_i = 0$) contribute zero columns to $\mat{B}_c$.
If all contacts are zero (fallback), all feet are treated as stance.

\textbf{Critical sign convention:} The GRF $\vect{f}_i$ is the force the
\emph{ground exerts on the foot} (positive $z$ = upward).  The skew
cross-product $\vect{r}_i \times \vect{f}_i$ already has the correct sign for
Newton-Euler Eq.~\eqref{eq:euler}.

\textbf{Frame convention:} Moment arms $\vect{r}_i$ are kept in the
\emph{world frame}.  The $\mat{B}$ matrix maps world-frame forces to
world-frame accelerations; mixing frames (e.g.\ rotating $\vect{r}$ into the
support plane while $\vect{f}$ remains in world frame) produces a physically
wrong cross product.  The support-plane rotation is used \emph{only} in the
friction-cone constraints.

\section{Zero-Order Hold Discretisation}

Given the continuous pair $(\mat{A}_c, \mat{B}_c)$ and a timestep $\Delta t$,
the ZOH discrete matrices are obtained from the matrix exponential:

\begin{equation}
    \begin{pmatrix}\mat{A}_d \\ \mat{B}_d\end{pmatrix}
    = \exp\!\left(
        \begin{pmatrix} \mat{A}_c & \mat{B}_c \\ \mat{0} & \mat{0} \end{pmatrix}
        \Delta t
    \right)_{:n_s,\,:}
    \label{eq:zoh}
\end{equation}

In practice, we form the $(n_s + n_u)\times(n_s + n_u)$ matrix
$\mat{M} = [\mat{A}_c, \mat{B}_c;\, \mat{0}, \mat{0}]\,\Delta t$,
compute $e^{\mat{M}}$ via \texttt{scipy.linalg.expm}, and read off:
\begin{equation}
    \mat{A}_d = e^{\mat{M}}[0\!:\!n_s,\; 0\!:\!n_s], \qquad
    \mat{B}_d = e^{\mat{M}}[0\!:\!n_s,\; n_s\!:]
\end{equation}

This is re-computed at \emph{every} MPC solve, because $\mat{B}_c$ changes
with foot positions and the contact schedule, making the system
\emph{time-varying} over the horizon.

% ══════════════════════════════════════════════════════════════════════════════
\chapter{Trot Gait Scheduler}
% ══════════════════════════════════════════════════════════════════════════════

\section{Diagonal Trot Pattern}

A trot gait alternates two \emph{diagonal pairs}:
\begin{itemize}[noitemsep]
  \item Pair 1: FL (front-left) + RR (rear-right)
  \item Pair 2: FR (front-right) + RL (rear-left)
\end{itemize}

\noindent\textbf{Parameters:}
\begin{align*}
    T &= 0.40\,\text{s}  \quad &&\text{(full gait cycle period, 2.5 Hz)}\\
    d &= 0.50            \quad &&\text{(duty cycle: 50\% stance, 50\% swing)}\\
    T_\text{stance} &= 0.20\,\text{s}, \quad T_\text{swing} = 0.20\,\text{s}
\end{align*}

\section{Phase Computation}

The phase offsets for each leg are:
\begin{equation}
    \text{offset} = [0.0,\; 0.5,\; 0.5,\; 0.0]
    \quad\text{for FL, FR, RL, RR respectively}
\end{equation}

Phase of leg $i$ at simulation time $t$:
\begin{equation}
    \phi_i(t) = \left(\frac{t - t_0}{T} + \mathrm{offset}_i\right) \bmod 1
    \label{eq:gait_phase}
\end{equation}

where $t_0$ is the gait \emph{activation time} (time when Phase C begins).
Leg $i$ is in stance when $\phi_i < d$, swing otherwise.

\textbf{Why relative phase matters (Bug Fix):}  If $t_0 = 0$ and the
simulation has been running for $t = 4.5\,\text{s}$, the phase
$4.5/0.4 = 11.25 \Rightarrow \phi = 0.25$ places the robot mid-cycle, causing
the first swing step to be incorrectly timed.  By anchoring to $t_0$ at
activation time, the gait always starts cleanly at $\phi = 0$.

\section{Horizon Schedule}

For the MPC, a $(K\times4)$ boolean schedule matrix is precomputed:

\begin{equation}
    \mat{S}[k, i] = \mathbf{1}\!\left[\phi_i(t_0^{\text{mpc}} + k\,\Delta t) < d\right]
    \quad k=0,\ldots,K-1
\end{equation}

where $t_0^{\text{mpc}}$ is the current simulation time at each MPC call.

% ══════════════════════════════════════════════════════════════════════════════
\chapter{State Estimation}
% ══════════════════════════════════════════════════════════════════════════════

\section{Ground-Truth Extraction}

Since this is a simulation-based controller, the state estimator reads
ground-truth values directly from MuJoCo data buffers every physics step
($2\,\text{ms}$ at $500\,\text{Hz}$):

\begin{itemize}[noitemsep]
  \item $\vect{p} = \texttt{d.qpos}[0\!:\!3]$ — CoM position.
  \item $\hat{\vect{q}} = \texttt{d.qpos}[3\!:\!7]$ (wxyz) — body quaternion,
        converted to rotation matrix $\mat{R}$ and Euler angles
        $[\phi, \theta, \psi]$ via ZYX convention.
  \item $\vect{v} = \texttt{d.qvel}[0\!:\!3]$ — CoM linear velocity (world).
  \item $\vect{\omega} = \texttt{d.qvel}[3\!:\!6]$ — angular velocity (world).
  \item Joint positions $\vect{q}$ and velocities $\dot{\vect{q}}$ read from
        their respective qpos/qvel slices.
  \item Foot world positions from \texttt{d.xpos[FOOT\_BODY\_IDS[leg]]}.
\end{itemize}

\section{Contact Detection Strategy}

\subsection{Raw efc\_force Reading}

MuJoCo reports constraint forces via \texttt{d.efc\_force[addr]}.  For each
active contact, the geom IDs are checked against a runtime-built map
(constructed from \texttt{mj\_name2id} lookups), and the force is accumulated
per leg:

\begin{equation}
    c_{f,i}^{\text{raw}} = \sum_{\text{contacts with foot }i}
    \left|\texttt{efc\_force}[\texttt{addr}]\right|
\end{equation}

\subsection{Rolling-Maximum Buffer}

A 10-step ($20\,\text{ms}$) rolling-maximum buffer suppresses single-step
dropouts:
\begin{equation}
    c_{f,i} = \max_{j \in \{0,\ldots,9\}}\; c_{f,i}^{\text{raw}}[t - j\,\Delta t_\text{phys}]
\end{equation}

\subsection{Height-Based Override (Critical Fix)}

MuJoCo's contact solver re-initialises its warm-start when the actuator
command vector changes abruptly, causing \texttt{efc\_force} to drop to
\emph{exactly zero} for up to $300\,\text{ms}$ even when all feet are
physically on the ground.  A rolling buffer of any practical size cannot
bridge this.

The correct fix uses the body height $p_z$ as a proxy:

\begin{equation}
    \text{contact}_i =
    \begin{cases}
        \texttt{True}  & \text{if } p_z > p_{z,\text{collapse}} = 0.20\,\text{m} \\
        c_{f,i} > \tau_\text{cf} & \text{otherwise}
    \end{cases}
    \label{eq:contact_override}
\end{equation}

with $\tau_\text{cf} = 10\,\text{N}$.  When $p_z > 0.20\,\text{m}$ the robot
is clearly standing; forcing all contacts true is always correct because:
\begin{enumerate}[noitemsep]
  \item Phase B uses all-stance regardless of contact state.
  \item Phase C uses the \emph{gait scheduler} for contact decisions,
        not this estimator.
  \item Collapse detection uses $p_z$ itself, so a false ``not in contact''
        cannot trigger it.
\end{enumerate}

% ══════════════════════════════════════════════════════════════════════════════
\chapter{Condensed Convex QP Formulation}
% ══════════════════════════════════════════════════════════════════════════════

\section{Prediction Horizon}

With MPC timestep $\Delta t_\text{mpc} = 30\,\text{ms}$ and horizon $K = 10$
steps, the controller looks $300\,\text{ms}$ ahead at each solve.

\section{Condensed Dynamics}

Define the stacked state trajectory $\mat{X} \in \R^{Kn_s}$ and stacked input
trajectory $\mat{U} \in \R^{Kn_u}$.  The time-varying discrete system
$\vect{x}_{k+1} = \mat{A}_k\vect{x}_k + \mat{B}_k\vect{u}_k$ gives the
\emph{condensed} form:

\begin{equation}
    \mat{X} = \mat{A}_{\mathrm{qp}}\,\vect{x}_0 + \mat{B}_{\mathrm{qp}}\,\mat{U}
    \label{eq:condensed}
\end{equation}

\subsection{Aqp Matrix}

$\mat{A}_{\mathrm{qp}} \in \R^{Kn_s \times n_s}$ stacks the state-transition
products:

\begin{equation}
    \mat{A}_{\mathrm{qp}} =
    \begin{pmatrix}
        \mat{A}_0 \\
        \mat{A}_1 \mat{A}_0 \\
        \vdots \\
        \mat{A}_{K-1}\cdots\mat{A}_0
    \end{pmatrix}
    \quad\text{where}\quad
    \Phi_k = \prod_{j=0}^{k}\mat{A}_j
\end{equation}

\subsection{Bqp Matrix}

$\mat{B}_{\mathrm{qp}} \in \R^{Kn_s \times Kn_u}$ is lower-block-triangular
(causal):

\begin{equation}
    \mat{B}_{\mathrm{qp}}[r,c] =
    \begin{cases}
        \Phi_{r,c}\,\mat{B}_c & \text{if } r \geq c \\
        \mat{0} & \text{otherwise}
    \end{cases}
    \quad\text{where}\quad
    \Phi_{r,c} = \prod_{j=c+1}^{r}\mat{A}_j
\end{equation}

\section{Cost Function}

\subsection{State Tracking}

Define the diagonal weight matrix $\mat{L} = \mat{I}_K \otimes \diag(\vect{q})$
where $\vect{q} \in \R^{13}$ are the state weights.  The state tracking cost is:

\begin{equation}
    J_\text{state} = \|\mat{X} - \mat{X}_\text{ref}\|^2_{\mat{L}}
    = (\mat{A}_{\mathrm{qp}}\vect{x}_0 + \mat{B}_{\mathrm{qp}}\mat{U} - \mat{X}_\text{ref})^\top
      \mat{L}
      (\mat{A}_{\mathrm{qp}}\vect{x}_0 + \mat{B}_{\mathrm{qp}}\mat{U} - \mat{X}_\text{ref})
\end{equation}

\subsection{Force Regularisation}

To prevent unrealistically large forces and ensure uniqueness when multiple
solutions exist:

\begin{equation}
    J_\text{force} = \alpha\,\|\mat{U}\|^2_2, \qquad \alpha = 3\times10^{-4}
\end{equation}

\subsection{Full QP}

\begin{equation}
    \min_{\mat{U}}\quad \frac{1}{2}\,\mat{U}^\top\mat{H}\,\mat{U} + \vect{g}^\top\mat{U}
    \label{eq:qp_cost}
\end{equation}

\begin{align}
    \mat{H} &= 2\left(\mat{B}_{\mathrm{qp}}^\top\mat{L}\,\mat{B}_{\mathrm{qp}}
               + \alpha\,\mat{I}_{Kn_u}\right) \\
    \vect{g} &= 2\,\mat{B}_{\mathrm{qp}}^\top\mat{L}
               \left(\mat{A}_{\mathrm{qp}}\,\vect{x}_0 - \mat{X}_\text{ref}\right)
\end{align}

\section{MPC Cost Weights}

\subsection{Flat-Ground Weights}

\begin{equation}
    \vect{q}_\text{flat} = \diag(10,\; 20,\; 1,\;\; 1,\; 1,\; 150,\;\; 1,\; 5,\; 1,\;\; 2,\; 2,\; 5,\;\; 0)
\end{equation}

\begin{center}
for $[\phi, \theta, \psi,\; p_x, p_y, p_z,\; \omega_x, \omega_y, \omega_z,\; v_x, v_y, v_z,\; {-g}]$
\end{center}

The high $p_z = 150$ strongly penalises height deviation, keeping the CoM at
the settled height.  The pitch weight $\theta = 20$ prevents forward/backward
tipping during trot.

\subsection{Stair-Climbing Weights}

\begin{equation}
    \vect{q}_\text{stairs} = \diag(60,\; 120,\; 1,\;\; 1,\; 15,\; 100,\;\; 8,\; 15,\; 1,\;\; 25,\; 15,\; 8,\;\; 0)
\end{equation}

The rationale for each change from the flat weights:

\begin{table}[H]
\centering
\caption{Stair Q-weight tuning rationale (4 iterations of log analysis)}
\begin{tabular}{p{1cm}p{2cm}p{2.5cm}p{7cm}}
\toprule
\textbf{State} & \textbf{Flat} & \textbf{Stairs} & \textbf{Rationale} \\
\midrule
$\phi$ (roll)   & 10  & 60  & Roll is the primary failure mode. Diagonal stance (FL on step, RR on flat) creates $\approx$9.3\,Nm roll torque. Equal priority with $p_z$. \\
$\theta$ (pitch)& 20  & 120 & Pitch diverges on stairs. Weight \emph{dominant} over $p_z$ (2:1) prevents height-tracking force from creating nose-up moment. \\
$p_y$           & 1   & 15  & Zero lateral authority caused $v_y$ to explode to 0.38\,m/s on first step. Strong lateral position tracking required. \\
$p_z$           & 150 & 100 & Reduced from flat (height tracking slightly sacrificed for pitch priority). Still 67\% of flat strength. \\
$\omega_x$      & 1   & 8   & Roll rate reached 15$^\circ$/s before divergence; strong damping needed. \\
$\omega_y$      & 5   & 15  & Pitch rate reached 10$^\circ$/s sustained; penalty 15$\times$0.17$^2$=0.44 provides meaningful deceleration. \\
$v_x$           & 2   & 25  & Robot surged 0.05$\to$0.21\,m/s (4$\times$ overspeed), leaving $<$1\,s per tread. Strong forward-speed regulation. \\
$v_y$           & 2   & 15  & Lateral velocity regulation to complement $p_y$ tracking. \\
$v_z$           & 5   & 8   & Vertical velocity oscillation ($\pm$0.15\,m/s) couples into pitch. \\
\bottomrule
\end{tabular}
\end{table}

\section{Friction Cone Constraints}

\subsection{Physical Constraint}

The foot can only push, not pull, and must stay within the friction cone:
\begin{equation}
    f_{n,i} \geq 0, \qquad \sqrt{f_{t,i}^2} \leq \mu\,f_{n,i}
\end{equation}

where $f_{n,i}$ is the normal force and $f_{t,i}$ is the tangential force,
$\mu = 0.7$ is the friction coefficient.

\subsection{Linearised Pyramid Approximation}

The circular cone is replaced by a square pyramid (4 planes):
\begin{align}
    f_{n,i} &\geq f_\text{min} = 5\,\text{N}
    \quad\text{(positive normal, min support force)} \\
    f_{n,i} &\leq f_\text{max} = 200\,\text{N}
    \quad\text{(leg load limit)} \\
    f_{t_x,i} - \mu\,f_{n,i} &\leq 0 \label{eq:fc1} \\
    -f_{t_x,i} - \mu\,f_{n,i} &\leq 0 \label{eq:fc2} \\
    f_{t_y,i} - \mu\,f_{n,i} &\leq 0 \label{eq:fc3} \\
    -f_{t_y,i} - \mu\,f_{n,i} &\leq 0 \label{eq:fc4}
\end{align}

Equations~\eqref{eq:fc1}--\eqref{eq:fc4} give 5 constraints per stance leg
($1$ normal bound $+$ 4 tangential planes).

\subsection{Terrain-Aware Rotation}

On sloped terrain (stairs), the support-plane normal $\hat{\vect{n}}$ deviates
from $[0,0,1]^\top$.  The friction cone is rotated into the support frame:

\begin{equation}
    \vect{f}_\text{plane} = \mat{R}_\text{plane}^\top\,\vect{f}_\text{world}
\end{equation}

where $\mat{R}_\text{plane}$ is the orthonormal basis of the support plane
(columns: $\hat{\vect{x}}_\text{plane}$, $\hat{\vect{y}}_\text{plane}$,
$\hat{\vect{n}}$).  The normal constraint then becomes:

\begin{equation}
    \hat{\vect{n}}^\top \vect{f}_i \in [f_\text{min},\; f_\text{max}]
\end{equation}

and the tangential constraints read:
\begin{equation}
    \hat{\vect{t}}_{x/y}^\top \vect{f}_i - \mu\,\hat{\vect{n}}^\top \vect{f}_i \leq 0
\end{equation}

\subsection{Swing-Leg Constraints}

For each swing leg $i$ ($S[k,i] = 0$), all three force components are zeroed:
\begin{equation}
    f_{x,i} = f_{y,i} = f_{z,i} = 0 \quad\text{(swing, 3 equality constraints)}
\end{equation}

% ══════════════════════════════════════════════════════════════════════════════
\chapter{Reference Trajectory Construction}
% ══════════════════════════════════════════════════════════════════════════════

\section{Operator Commands}

The operator specifies body-frame commands:
\begin{equation}
    \vect{c} = [v_x^\text{cmd},\; v_y^\text{cmd},\; \omega_z^\text{cmd}]
\end{equation}

\section{Per-Step Reference}

For each horizon step $k$ at projected position $x_k = p_x + v_x^\text{cmd}\,k\,\Delta t$:

\begin{equation}
    \vect{r}_k =
    \begin{pmatrix}
        0                          \\ % roll
        \theta_\text{ref}          \\ % pitch
        \psi_0 + \omega_z^\text{cmd}\,k\,\Delta t \\ % yaw
        x_k                        \\ % px
        p_y + v_y^\text{cmd}\,k\,\Delta t \\ % py
        p_{z,\text{des}}^{(k)}     \\ % pz
        0                          \\ % ωx
        0                          \\ % ωy
        \omega_z^\text{cmd}        \\ % ωz
        v_x^\text{cmd}             \\ % vx
        v_y^\text{cmd}             \\ % vy
        0                          \\ % vz
        -g                         % gravity augmentation
    \end{pmatrix}
    \label{eq:Xref}
\end{equation}

\section{Terrain-Aware Reference (Stairs)}

When terrain mode is active, both $\theta_\text{ref}$ and $p_{z,\text{des}}$
are computed as a \emph{delta from the current projected position}:

\begin{align}
    \theta_\text{ref}^{(k)} &= \theta_\text{ref,smooth} + \Delta\theta(x_k) \\
    \Delta\theta(x_k)       &= \texttt{pitch\_ref}(x_k) - \texttt{pitch\_ref}(x_0)
\end{align}

\begin{align}
    p_{z,\text{des}}^{(k)} &= p_{z,\text{des,smooth}} + \Delta p_z(x_k) \\
    \Delta p_z(x_k) &= \texttt{target\_pz}(x_k, \ell) - \texttt{target\_pz}(x_0, \ell)
\end{align}

where $\ell$ is the nominal CoM clearance above terrain.  This delta
formulation preserves the IIR-smoothed current values at $k=0$ while
anticipating terrain changes in future horizon steps.

\section{Height Integral Correction}

A small integral correction prevents steady-state height error:
\begin{equation}
    p_{z,\text{corrected}} = p_{z,\text{des}} + \text{clip}\!\left(0.1\,e_{z,\text{int}},\;-0.05,\;0.05\right)
\end{equation}

% ══════════════════════════════════════════════════════════════════════════════
\chapter{Swing Leg Controller}
% ══════════════════════════════════════════════════════════════════════════════

\section{Minimum-Jerk Trajectory}

The swing trajectory interpolates from the lift-off position $\vect{p}_0$
to the touchdown position $\vect{p}_f$ using a \emph{minimum-jerk} polynomial:

\begin{equation}
    s(t) = \frac{t - t_\text{liftoff}}{T_\text{swing}}, \quad s \in [0, 1]
\end{equation}

\begin{equation}
    P_\text{mj}(s) = 10s^3 - 15s^4 + 6s^5
    \label{eq:minjerk}
\end{equation}

The position trajectory is:
\begin{equation}
    \vect{p}_\text{des}(s) = \vect{p}_0 + (\vect{p}_f - \vect{p}_0)\,P_\text{mj}(s)
                             + z_\text{bump}(s)\,\hat{\vect{e}}_z
\end{equation}

\subsection{Height Bump}

A smooth height bump ensures foot clearance above the terrain:

\begin{equation}
    z_\text{bump}(s) = h_\text{sw}\cdot 64\,s^3(1-s)^3
    \label{eq:heightbump}
\end{equation}

Properties:
\begin{itemize}[noitemsep]
  \item $z_\text{bump}(0) = z_\text{bump}(1) = 0$ (no displacement at endpoints).
  \item Peak at $s = 0.5$: $z_\text{bump}(0.5) = h_\text{sw}$.
  \item $\dot{z}_\text{bump}(0) = \dot{z}_\text{bump}(1) = 0$ (smooth lift/land).
\end{itemize}

For flat ground: $h_\text{sw} = 0.08\,\text{m}$.  On stairs: $h_\text{sw}$
is terrain-adaptive (see Chapter~\ref{ch:stairs}).

\subsection{Velocity and Acceleration}

\begin{align}
    \dot{P}_\text{mj}(s) &= 30s^2 - 60s^3 + 30s^4 \\
    \ddot{P}_\text{mj}(s) &= 60s - 180s^2 + 120s^3
\end{align}

Foot velocity: $\vect{v}_\text{des} = (\vect{p}_f - \vect{p}_0)\,\dot{P}_\text{mj}/T_\text{swing}$

\section{Raibert Touchdown Heuristic}

The foot touchdown position is determined by the Raibert heuristic
\cite{raibert1986}:

\begin{equation}
    \vect{p}_\text{td} = \vect{p}_\text{hip} +
    \underbrace{\vect{v}_\text{cmd}\,t_\text{pred}}_{\text{nominal drift}} +
    \underbrace{k_v\,(\vect{v}_\text{actual} - \vect{v}_\text{cmd})}_{\text{error correction}}
    \label{eq:raibert}
\end{equation}

where:
\begin{align}
    t_\text{pred} &= \frac{T}{2} = \frac{T_\text{swing} + 0.5\,T_\text{stance}}{2}
    \approx 0.20\,\text{s} \\
    k_{v,x} &= 0.4\,T \approx 0.12\,\text{m\,s} \\
    k_{v,y} &= 0.05\,T \approx 0.02\,\text{m\,s}  \quad\text{(reduced to prevent lateral instability)}
\end{align}

\textbf{Command frame:} The commanded velocity is in the body frame.  It is
rotated to the world frame before applying the Raibert formula:
\begin{equation}
    \begin{pmatrix}v_x^w \\ v_y^w\end{pmatrix} =
    \begin{pmatrix}\cos\psi & -\sin\psi \\ \sin\psi & \cos\psi\end{pmatrix}
    \begin{pmatrix}v_x^\text{cmd} \\ v_y^\text{cmd}\end{pmatrix}
\end{equation}

\textbf{Overcorrection fix:} The original implementation used
$\vect{p}_\text{td} = \vect{p}_\text{hip} + \vect{v}_\text{actual}\,(t_\text{pred} + k_v)$,
which applies $(t_\text{pred} + k_v)\,v_\text{actual}$ as one block.  At a
$0.51\,\text{m/s}$ lateral drift with $k_v = 0.08$, this gives a 3.5$\times$
larger correction than intended.  The corrected formulation separates nominal
drift (proportional to \emph{commanded} velocity) from error correction
(proportional to velocity \emph{error}).

\section{Task-Space PD Controller}

For each swing leg, the joint torques are computed as:

\begin{equation}
    \vect{\tau}_\text{swing} = \mat{J}_i^\top \underbrace{\left[
        \mat{K}_p\,(\vect{p}_\text{des} - \vect{p}_\text{cur}) +
        \mat{K}_d\,(\vect{v}_\text{des} - \vect{v}_\text{cur})
    \right]}_{\vect{F}_\text{task}} + \vect{\tau}_\text{bias}
    \label{eq:swing_pd}
\end{equation}

where:
\begin{align}
    \mat{K}_p &= \diag(400, 400, 400)\,\text{N/m} \\
    \mat{K}_d &= \diag(75, 75, 75)\,\text{N\,s/m} \\
    \mat{J}_i &= \text{3×3 foot Jacobian (leg DOFs)} \\
    \vect{\tau}_\text{bias} &= \texttt{qfrc\_bias}[\text{leg joints}]
    \quad\text{(gravity/Coriolis compensation)}
\end{align}

\textbf{Foot velocity:} The full foot velocity uses the complete
$(3\times n_v)$ Jacobian (body + all leg DOFs):
\begin{equation}
    \vect{v}_\text{cur} = \mat{J}_\text{full}\;\dot{\vect{q}}_\text{full}
\end{equation}
This includes body translational/angular velocity, which is significant when
the body is rolling or pitching.

\section{GRF to Joint Torques (Stance)}

For stance legs, the MPC ground-reaction forces are converted to joint torques:

\begin{equation}
    \vect{\tau}_\text{stance,i} = -\mat{J}_i^\top\;\vect{f}_i + \vect{\tau}_\text{bias,i}
    \label{eq:stance_tau}
\end{equation}

\textbf{Sign convention (critical):} The MPC gives $\vect{f}_i$ as the force
the \emph{ground exerts on the foot} (upward positive).  To generate this
upward reaction, the joint must push \emph{down} against the ground.
For a knee joint with $J_z < 0$ at nominal stance, $-\mat{J}^\top\vect{f}$
gives extension torque (pushing down), which is physically correct.

\section{Actuator Reordering}

MuJoCo's actuator order for the Go2 differs from the internal leg ordering:

\begin{center}
\begin{tabular}{lll}
\toprule
\textbf{Internal order} & \textbf{Indices} & \textbf{MuJoCo ctrl indices} \\
\midrule
FL & 0--2  & ctrl[3--5] \\
FR & 3--5  & ctrl[0--2] \\
RL & 6--8  & ctrl[9--11] \\
RR & 9--11 & ctrl[6--8] \\
\bottomrule
\end{tabular}
\end{center}

The \texttt{leg\_to\_ctrl()} function performs this remapping before writing
to \texttt{data.ctrl}.

% ══════════════════════════════════════════════════════════════════════════════
\chapter{Controller State Machine}
% ══════════════════════════════════════════════════════════════════════════════

\section{Three-Phase Architecture}

\begin{figure}[H]
\centering
\begin{tikzpicture}[
    node distance=3.5cm,
    state/.style={rectangle, rounded corners, draw, fill=blue!10,
                  minimum width=3cm, minimum height=1.2cm, align=center},
    arrow/.style={->, thick}
]
\node[state] (A) {Phase A\\PD Stand\\$0\!-\!1.5\,\text{s}$};
\node[state, right=of A] (B) {Phase B\\All-stance MPC\\$1.5\!-\!3.5\,\text{s}$};
\node[state, right=of B] (C) {Phase C\\Trot MPC\\$t > 3.5\,\text{s}$};
\node[state, below=1.5cm of B, fill=red!10] (R) {Recovery\\PD Stand\\2\,s};

\draw[arrow] (A) -- node[above, font=\small]{$t > 1.5\,\text{s}$} (B);
\draw[arrow] (B) -- node[above, font=\small]{$t > 3.5\,\text{s}$, $p_z\geq0.30\,\text{m}$} (C);
\draw[arrow] (C) to[bend left] node[right, font=\small]{$p_z < 0.20\,\text{m}$} (R);
\draw[arrow] (R) to[bend left] node[left, font=\small]{$+2\,\text{s}$} (B);
\end{tikzpicture}
\caption{MPC controller three-phase state machine}
\end{figure}

\section{Phase A: PD Settling ($0$--$1.5\,\text{s}$)}

\subsection{Goal}
Bring the robot from its initial (possibly sprawled) configuration to a
nominal standing posture with stable contact on all four feet.

\subsection{Controller}
A joint-space PD controller tracks nominal standing angles:

\begin{equation}
    \vect{\tau}_\text{PD} = \mat{K}_P\,(\vect{q}_\text{des} - \vect{q}) -
    \mat{K}_D\,\dot{\vect{q}}
\end{equation}

with $\vect{q}_\text{des} = [0, 0.7, -1.4]$ rad per leg and
$\mat{K}_P = \mat{K}_D = \diag(80, 80, 80)\,\text{N\,m/rad}$,
$\mat{K}_D = \diag(8, 8, 8)\,\text{N\,m\,s/rad}$.

\subsection{Height Estimation}
During the last $0.8\,\text{s}$ of Phase A, an IIR filter tracks the settled
CoM height:
\begin{equation}
    p_{z,\text{settled}} \leftarrow 0.98\,p_{z,\text{settled}} + 0.02\,p_z
\end{equation}
This settled height ($\approx 0.338\,\text{m}$) is used as $p_{z,\text{des}}$
for all subsequent phases.

\section{Phase B: All-Stance MPC ($1.5$--$3.5\,\text{s}$)}

\subsection{Goal}
Smoothly transition from PD joint-space control to GRF-based force control
while maintaining full 4-foot contact.  This ``warm-starts'' the MPC with
a tractable (all-stance) problem before introducing the complexity of leg swing.

\subsection{Contact Schedule}
All 4 feet in stance for the full horizon:
\begin{equation}
    \mat{S}[k, i] = 1 \quad \forall k, i
\end{equation}

\subsection{GRF Initialisation}
Seeded at the static equilibrium:
\begin{equation}
    f_{z,i}^{(0)} = \frac{mg}{4} = \frac{15.206 \times 9.81}{4} \approx 37.3\,\text{N}
\end{equation}

\subsection{Phase B Warmup (100 ms)}
During the first $100\,\text{ms}$ of Phase B, the MPC runs at the physics
step rate ($2\,\text{ms}$) to rapidly build correct GRF estimates from the
fresh measurement state.  After warmup, the rate drops to the nominal
$30\,\text{ms}$ cadence.

\section{Phase C: Trot MPC ($t > 3.5\,\text{s}$)}

\subsection{Entry Gate}
Phase C is entered only when:
\begin{enumerate}[noitemsep]
  \item $t > 3.5\,\text{s}$ (Phase B duration elapsed).
  \item $p_z \geq 0.30\,\text{m}$ (robot is standing tall enough).
  \item $|\theta| < 2^\circ$ (body is level — prevents trot on tilted stance).
\end{enumerate}

\subsection{Phase C Warmup (300 ms)}
After Phase C is entered, the gait scheduler is activated but all 4 feet
remain in stance for $300\,\text{ms}$ (1.5 gait cycles).  This gives the
MPC time to smooth forces before the first swing leg lifts off.

\subsection{Normal Operation}
After warmup:
\begin{enumerate}[noitemsep]
  \item Gait scheduler provides the $(K\times4)$ contact schedule.
  \item Lift-off transitions trigger \texttt{notify\_lift()} on the swing
        controller (Raibert touchdown computed, trajectory initialised).
  \item Stance legs: torques from MPC GRFs via $-\mat{J}^\top\vect{f}$.
  \item Swing legs: torques from task-space PD via \eqref{eq:swing_pd}.
  \item GRF solution is smoothed before application (Section~\ref{sec:grf_smooth}).
\end{enumerate}

\section{GRF Smoothing}
\label{sec:grf_smooth}

A first-order IIR filter blends consecutive QP solutions:

\begin{equation}
    \vect{f}^{(t)} = \beta\,\vect{f}_\text{new} + (1 - \beta)\,\vect{f}^{(t-1)}
\end{equation}

\begin{align}
    \beta_\text{flat}  &= 0.75 \\
    \beta_\text{stairs}&= 0.65
\end{align}

On stairs, a slower $\beta$ prevents force spikes when contact geometry
changes abruptly.

\subsection{GRF Re-Seeding}

When a leg transitions from swing to stance, its stored GRF is near zero.
If it falls below 70\% of the static equilibrium value, it is immediately
re-seeded:

\begin{equation}
    \text{if } f_{z,i} < 0.7\,\frac{mg}{n_\text{stance}} \quad
    \Rightarrow\quad \vect{f}_i \leftarrow \left[0,\;0,\;\frac{mg}{n_\text{stance}}\right]
\end{equation}

This eliminates the underforce window ($\approx$100\,ms) that occurs when
the smoothed GRF is too low on the first tick of stance.

\section{Collapse Detection and Recovery}

If $p_z < 0.20\,\text{m}$ while the MPC is active:
\begin{enumerate}[noitemsep]
  \item Switch to PD stand controller.
  \item Hold PD for $2.0\,\text{s}$.
  \item Return to Phase B (all-stance MPC).
\end{enumerate}

% ══════════════════════════════════════════════════════════════════════════════
\chapter{Flat-Ground Trotting: End-to-End Flow}
% ══════════════════════════════════════════════════════════════════════════════

\section{Step-by-Step Control Loop}

Each physics step ($\Delta t_\text{phys} = 2\,\text{ms}$):

\begin{enumerate}
    \item \textbf{State read.} \texttt{StateEstimator.update()} reads
          $\vect{x}$, joint positions, foot positions, contact forces.

    \item \textbf{Contact detection.} Rolling-max buffer + height override
          (Eq.~\ref{eq:contact_override}).

    \item \textbf{Height integral.}
          $e_{z,\text{int}} \mathrel{+}= (p_{z,\text{des}} - p_z)\,\Delta t_\text{phys}$.

    \item \textbf{Phase logic.} Select Phase A/B/C based on simulation time
          and height.

    \item \textbf{MPC solve} (every $30\,\text{ms}$, rate-limited):
        \begin{enumerate}
            \item Build $\vect{x}_0 = \texttt{state.mpc\_x()}$.
            \item Compute gait schedule $\mat{S} \in \{0,1\}^{K\times4}$.
            \item Build $\mat{A}_c, \mat{B}_c$ for each horizon step $k$.
            \item Discretise: $(\mat{A}_d^{(k)}, \mat{B}_d^{(k)}) =
                  \text{ZOH}(\mat{A}_c^{(k)}, \mat{B}_c^{(k)}, \Delta t_\text{mpc})$.
            \item Condense: $(\mat{A}_\text{qp}, \mat{B}_\text{qp})$.
            \item Build reference $\mat{X}_\text{ref}$.
            \item Compute $\mat{H}, \vect{g}$ (cost).
            \item Build constraints $\mat{C}, \vect{lb}, \vect{ub}$.
            \item Solve QP $\Rightarrow \mat{U}_\text{opt}$.
            \item Extract $\vect{f}_{1:4} = \mat{U}_\text{opt}[0\!:\!12]$.
            \item Apply GRF smoothing.
        \end{enumerate}

    \item \textbf{Swing detection.} Check lift-off transitions
          ($\phi_i$ crossed duty boundary); call \texttt{notify\_lift()} with
          Raibert touchdown target.

    \item \textbf{Torque computation.}
          \begin{itemize}[noitemsep]
            \item Stance: $\vect{\tau}_i = -\mat{J}_i^\top\vect{f}_i + \vect{\tau}_\text{bias,i}$.
            \item Swing: $\vect{\tau}_i = \mat{J}_i^\top\vect{F}_\text{task,i} + \vect{\tau}_\text{bias,i}$.
          \end{itemize}

    \item \textbf{Actuator remapping.} Internal $\to$ MuJoCo order.

    \item \textbf{Torque limit.} Clip to $\pm 33.5\,\text{N\,m}$ per joint.

    \item \textbf{Apply.} \texttt{data.ctrl[:] = ctrl}; \texttt{mj\_step()}.

    \item \textbf{Log.} Write row to CSV debug log.
\end{enumerate}

% ══════════════════════════════════════════════════════════════════════════════
\chapter{Stair Climbing System}
\label{ch:stairs}
% ══════════════════════════════════════════════════════════════════════════════

\section{Stair Scene Geometry}

The \texttt{scene\_stairs\_uniform.xml} scene contains:

\begin{table}[H]
\centering
\caption{Stair scene geometry}
\begin{tabular}{llll}
\toprule
\textbf{Section} & \textbf{x range} & \textbf{Surface $z$} & \textbf{Steps} \\
\midrule
Flat approach   & $x < 2.00\,\text{m}$           & $0.00\,\text{m}$  & — \\
Ascending stairs& $2.00 \leq x < 3.40\,\text{m}$ & $0.06 \to 0.30\,\text{m}$ & 5 \\
Top platform    & $3.40 \leq x < 4.40\,\text{m}$ & $0.30\,\text{m}$  & — \\
Descending stairs & $4.40 \leq x < 5.80\,\text{m}$ & $0.24 \to 0.06\,\text{m}$ & 5 \\
Flat exit       & $x \geq 5.80\,\text{m}$         & $0.00\,\text{m}$  & — \\
\bottomrule
\end{tabular}
\end{table}

\noindent\textbf{Step parameters:}
\begin{align*}
    &\text{Riser height:} \quad h_r = 0.06\,\text{m} \\
    &\text{Tread width:}  \quad w_t = 0.28\,\text{m} \\
    &\text{Total height:} \quad H = 5 \times 0.06 = 0.30\,\text{m}
\end{align*}

\section{Terrain Mode Activation}

The controller initialises in ``flat'' mode (terrain estimator is a no-op).
The simulation loop calls \texttt{switch\_to\_terrain\_mode(``stairs'')} when
the robot's $x$-position exceeds the approach threshold
($x > 1.80\,\text{m}$, i.e.\ $0.20\,\text{m}$ before the first riser).

At the switch:
\begin{enumerate}[noitemsep]
  \item Terrain estimator mode set to ``stairs''.
  \item Nominal clearance anchored: $\ell = p_{z,\text{settled}} - z_\text{terrain}$.
  \item Stairs-specific Q weights activated immediately.
  \item Swing clearance changes: $0.08\,\text{m} \to 0.18\,\text{m}$ (adaptive).
\end{enumerate}

\section{Terrain Height Map}

\subsection{Discrete Step Function}

Used for touchdown z initialisation and failure detection:

\begin{equation}
    z_\text{terrain}(x) =
    \begin{cases}
        0 & x < 2.00 \\
        \lceil\frac{x - 2.00}{0.28}\rceil \times 0.06 & 2.00 \leq x < 3.40
                \quad\text{(ascending)} \\
        0.30 & 3.40 \leq x < 4.40 \quad\text{(platform)} \\
        0.30 - \lceil\frac{x - 4.40}{0.28}\rceil \times 0.06 & 4.40 \leq x < 5.80
                \quad\text{(descending)} \\
        0 & x \geq 5.80
    \end{cases}
\end{equation}

\subsection{Smooth Ramp Approximation}

Used for pitch\_ref and target\_pz to ensure smooth reference transitions:

An $8\,\text{cm}$ linear pre-ramp before each riser edge blends the height
continuously, preventing step-changes in the MPC reference.

\section{Adaptive CoM Height}

\subsection{Target Height}

\begin{equation}
    p_{z,\text{target}}(x, \ell) =
    \underbrace{\frac{z_\text{smooth}(x + d_\text{hip}) + z_\text{smooth}(x - d_\text{hip})}{2}}_{\text{avg terrain under both hips}}
    + \ell
\end{equation}

where $d_\text{hip} = 0.1934\,\text{m}$ and $\ell \approx 0.338\,\text{m}$ (nominal clearance).

\subsection{Asymmetric Rate Limiter}

\begin{equation}
    \Delta p_{z,\text{des}} = \text{clip}\!\left(p_{z,\text{target}} - p_{z,\text{des}},\;
    -r_\text{down}\,\Delta t,\; r_\text{up}\,\Delta t\right)
\end{equation}

\begin{align}
    r_\text{up}   &= 0.12\,\text{m/s} \quad\text{(ascending — fast to prevent lag)} \\
    r_\text{down} &= 0.04\,\text{m/s} \quad\text{(descending — slow to prevent drops)}
\end{align}

This means a $6\,\text{cm}$ riser is tracked in $\approx 250\,\text{ms}$ on
ascent (vs.\ $1.5\,\text{s}$ at the original $0.04\,\text{m/s}$ symmetric rate
which caused sustained overforce and roll crash).

\section{Body Pitch Reference}

\subsection{Geometric Derivation}

The desired body pitch angle is computed from the terrain slope under the hip
span:

\begin{equation}
    \theta_\text{ref} = \text{clip}\!\left(
        \arctan\!\left(\frac{z_\text{smooth}(x + d_\text{hip}) - z_\text{smooth}(x - d_\text{hip})}
                             {2\,d_\text{hip}}\right),\;
        -0.21, 0.21\right)
    \label{eq:pitch_ref}
\end{equation}

\begin{itemize}[noitemsep]
  \item Ascending ($2.00 \to 3.40\,\text{m}$): $\theta_\text{ref} \approx +0.154\,\text{rad} \approx +8.8^\circ$ (nose-up).
  \item Platform ($3.40 \to 4.40\,\text{m}$): $\theta_\text{ref} = 0$.
  \item Descending ($4.40 \to 5.80\,\text{m}$): $\theta_\text{ref} \approx -0.154\,\text{rad} \approx -8.8^\circ$ (nose-down).
\end{itemize}

\textbf{Physical reason for negative pitch on descent:}  When front feet
are on lower steps and rear feet are still higher, a level body shifts the CoM
forward of the support polygon, creating a nose-dive torque.  By commanding
$-8.8^\circ$ pitch reference, the MPC proactively distributes forces to
maintain this geometry throughout descent.

\textbf{Flat approach guard:}  The smooth height function has an $8\,\text{cm}$
pre-ramp before $x = 2.00\,\text{m}$.  This would generate a phantom
$+8^\circ$ pitch command on flat ground.  The guard returns $\theta_\text{ref} = 0$
when $x < 2.00\,\text{m}$.

\subsection{IIR Smoothing}

The geometric pitch reference is filtered before passing to the MPC:

\begin{equation}
    \theta_\text{ref,smooth} \leftarrow 0.3\,\theta_\text{ref,smooth} + 0.7\,\theta_\text{ref,raw}
\end{equation}

The fast smoothing constant ($\alpha = 0.3$, settling in $\approx 2$ ticks)
prevents the MPC from generating anti-pitch torques during the $3$--$6^\circ$
lag that occurred with $\alpha = 0.5$.

\section{Per-Foot Stair Touchdown Placement}

\subsection{Tread Snapping}

The Raibert heuristic computes a continuous touchdown x-position.  On stairs,
this must be clamped to a safe zone within the target tread (margin: $4\,\text{cm}$
from each riser edge):

\begin{equation}
    x_\text{td,safe} = \text{clip}\!\left(x_\text{Raibert},\;
    x_\text{tread,lo} + \delta, \;
    x_\text{tread,hi} - \delta \right), \quad \delta = 0.04\,\text{m}
\end{equation}

\subsection{Height Map Touchdown}

The touchdown z-coordinate is set from the deterministic height map:

\begin{equation}
    z_\text{td} = z_\text{terrain}(x_\text{td,safe}) + r_\text{foot}
\end{equation}

where $r_\text{foot} = 0.022\,\text{m}$ is the foot sphere radius.

This ensures the foot targets the correct tread surface even if the body has
not yet risen to track the terrain.

\section{Adaptive Swing Clearance}

\begin{equation}
    h_\text{sw}(z_\text{terrain}) =
    h_\text{flat} + \min\!\left(1,\; \frac{z_\text{terrain}}{z_\text{ramp}}\right)
    \times (h_\text{terrain} - h_\text{flat} - 0.01) + 0.01
\end{equation}

\begin{align*}
    h_\text{flat} &= 0.08\,\text{m},\quad
    h_\text{terrain} = 0.18\,\text{m},\quad
    z_\text{ramp} = 0.10\,\text{m}
\end{align*}

As the terrain IIR estimate rises from $0 \to 0.30\,\text{m}$,
swing clearance grows from $0.08 \to 0.18\,\text{m}$, ensuring
the foot clears the step it is swinging to land on.

\section{Early Touchdown Detection}

When a swing leg physically contacts the terrain before its scheduled
touchdown (e.g.\ the foot catches the riser face), the controller must
stop pushing it forward and freeze it in place:

\begin{equation}
    \text{early\_td}[i] = \mathbf{1}\!\left[
        c_{f,i} > 10\,\text{N}
        \;\wedge\;
        \frac{t - t_\text{liftoff}}{T_\text{swing}} > 0.60
    \right]
    \label{eq:early_td}
\end{equation}

The two guards in Eq.~\eqref{eq:early_td} are essential:

\begin{enumerate}
  \item \textbf{Guard 1 (terrain mode only):} In flat mode, early touchdown
        is disabled.  The rolling-max buffer retains stance contact forces for
        $\approx 20\,\text{ms}$ post-liftoff, causing false triggers.

  \item \textbf{Guard 2 (60\% swing progress):} Feet hitting the \emph{riser
        face} at 20--30\% swing progress would freeze far short of the safe
        tread zone.  Requiring 60\% ensures the foot has crested the riser
        and is descending toward the tread before early touchdown fires.
\end{enumerate}

When early touchdown is active:
\begin{equation}
    \vect{\tau}_i = \mat{J}_i^\top\!\left[\mat{K}_d\,(-\vect{v}_\text{foot})\right]
    + \vect{\tau}_\text{bias,i}
    \quad\text{(velocity damping only, no position error)}
\end{equation}

\section{Support Plane Estimation}

\subsection{Purpose}

When stance feet are on different stair steps, they are not coplanar.
The support plane provides a best-fit terrain normal for rotating friction-cone
constraints.

\subsection{Least-Squares Plane Fit}

Given $n \geq 3$ stance foot positions $\{\vect{p}_i\}$:

\begin{enumerate}
  \item Compute centroid: $\vect{c} = \frac{1}{n}\sum_i \vect{p}_i$.
  \item Form $\mat{A} = [\vect{p}_1 - \vect{c},\; \ldots,\; \vect{p}_n - \vect{c}]^\top$.
  \item Compute SVD: $\mat{A} = \mat{U}\mat{\Sigma}\mat{V}^\top$.
  \item Normal: $\hat{\vect{n}} = \mat{V}[-1,:]$ (last row of $\mat{V}^\top$,
        smallest singular value direction).
  \item Flip if $n_z < 0$.
\end{enumerate}

\subsection{Support Frame Construction}

\begin{align}
    \hat{\vect{n}} &= \hat{\vect{n}} / \|\hat{\vect{n}}\| \\
    \hat{\vect{x}}_\text{plane} &= \hat{\vect{x}}_\text{world} -
    (\hat{\vect{x}}_\text{world}\cdot\hat{\vect{n}})\,\hat{\vect{n}}
    \quad\text{(project world-}x\text{ onto plane)} \\
    \hat{\vect{y}}_\text{plane} &= \hat{\vect{n}} \times \hat{\vect{x}}_\text{plane} \\
    \mat{R}_\text{plane} &= [\hat{\vect{x}}_\text{plane},\;
                              \hat{\vect{y}}_\text{plane},\;
                              \hat{\vect{n}}]
\end{align}

$\mat{R}_\text{plane}$ is passed to the QP constraints for friction-cone
rotation but \emph{not} used in the $\mat{B}$ matrix (moment arms stay in
world frame).

\section{IIR Terrain Tracking}

The terrain estimator maintains an IIR-filtered estimate of the terrain height
from stance foot z-positions:

\begin{equation}
    z_\text{terrain} \leftarrow 0.70\,z_\text{terrain} + 0.30\,\max\!\left(0,\;
    \overline{z}_\text{stance} - r_\text{foot}\right)
\end{equation}

where $\overline{z}_\text{stance}$ is the mean z-position of all feet with
$c_{f,i} > 10\,\text{N}$.  This estimate drives the adaptive swing clearance
and is used as the ``current terrain'' baseline for touchdown offset.

% ══════════════════════════════════════════════════════════════════════════════
\chapter{Complete Constraints Reference}
% ══════════════════════════════════════════════════════════════════════════════

\begin{longtable}{p{4cm}p{4cm}p{6cm}}
\caption{All constraints active in the QP and elsewhere}\\
\toprule
\textbf{Constraint} & \textbf{Equation} & \textbf{Purpose / Notes} \\
\midrule
\endfirsthead
\multicolumn{3}{l}{\textit{(continued)}}\\
\toprule
\textbf{Constraint} & \textbf{Equation} & \textbf{Purpose / Notes} \\
\midrule
\endhead
\bottomrule
\endfoot

Normal force min   & $\hat{\vect{n}}^\top\vect{f}_i \geq 5\,\text{N}$
                   & Prevent foot from losing contact; ensure minimum support \\
Normal force max   & $\hat{\vect{n}}^\top\vect{f}_i \leq 200\,\text{N}$
                   & Leg structural limit \\
Friction cone $+x$ & $\hat{\vect{t}}_x^\top\vect{f}_i \leq \mu\,\hat{\vect{n}}^\top\vect{f}_i$
                   & Linearised pyramid, positive tangential \\
Friction cone $-x$ & $-\hat{\vect{t}}_x^\top\vect{f}_i \leq \mu\,\hat{\vect{n}}^\top\vect{f}_i$
                   & Linearised pyramid, negative tangential \\
Friction cone $+y$ & $\hat{\vect{t}}_y^\top\vect{f}_i \leq \mu\,\hat{\vect{n}}^\top\vect{f}_i$
                   & Linearised pyramid, positive lateral \\
Friction cone $-y$ & $-\hat{\vect{t}}_y^\top\vect{f}_i \leq \mu\,\hat{\vect{n}}^\top\vect{f}_i$
                   & Linearised pyramid, negative lateral \\
Swing zero force   & $\vect{f}_i = \vect{0}$
                   & Swing leg cannot exert force (3 equality constraints) \\
Torque limit       & $|\tau_j| \leq 33.5\,\text{N\,m}$
                   & Hardware safety limit, applied after all torque computation \\
Collapse gate      & $p_z \geq 0.20\,\text{m}$
                   & Trigger PD recovery if violated \\
Trot gate          & $p_z \geq 0.30\,\text{m}$,\; $|\theta| < 2^\circ$
                   & Phase C entry condition \\
pz rate limit up   & $\dot{p}_{z,\text{des}} \leq 0.12\,\text{m/s}$
                   & Smooth height increase on stairs \\
pz rate limit down & $\dot{p}_{z,\text{des}} \geq -0.04\,\text{m/s}$
                   & Smooth height decrease on stairs \\
Contact threshold  & $c_{f,i} > 10\,\text{N}$
                   & Minimum force to classify as stance (for early touchdown and state estimator) \\
Early touchdown    & $\phi_\text{swing} > 0.60$
                   & Guard against premature early touchdown detection \\
Tread margin       & $x \in [x_\text{tread} + 0.04,\; x_\text{tread+1} - 0.04]$
                   & Safe foot placement within each step \\
Pitch clip         & $|\theta_\text{ref}| \leq 0.21\,\text{rad}$
                   & Cap pitch reference at $\approx 12^\circ$ \\
\end{longtable}

% ══════════════════════════════════════════════════════════════════════════════
\chapter{Module Dependency and Data Flow}
% ══════════════════════════════════════════════════════════════════════════════

\section{Module Dependency Graph}

\begin{figure}[H]
\centering
\begin{tikzpicture}[
    node distance=2cm,
    mod/.style={rectangle, draw, fill=blue!8, rounded corners,
                minimum width=3cm, minimum height=0.7cm, align=center, font=\small},
    arr/.style={->, gray!60, thick}
]
\node[mod] (cfg) {config.py};
\node[mod, below=0.5cm of cfg] (rp) {robot\_params.py};
\node[mod, right=3cm of cfg] (dyn) {dynamics.py};
\node[mod, right=3cm of rp] (sol) {mpc\_solver.py};
\node[mod, below=0.5cm of rp] (se) {state\_estimator.py};
\node[mod, below=0.5cm of se] (gs) {gait\_scheduler.py};
\node[mod, below=0.5cm of gs] (sw) {swing\_controller.py};
\node[mod, right=3cm of se] (tu) {torque\_utils.py};
\node[mod, right=3cm of gs] (te) {terrain\_estimator.py};
\node[mod, right=3cm of sw] (sp) {support\_plane.py};
\node[mod, below=3cm of sol] (mc) {mpc\_controller.py};
\node[mod, below=0.8cm of mc] (sim) {simulation.py};
\node[mod, below left=0.5cm and 0.5cm of sim] (ft) {flat\_trot.py};
\node[mod, below right=0.5cm and 0.5cm of sim] (sc) {stairs\_climb.py};

\draw[arr] (cfg) -- (dyn);
\draw[arr] (cfg) -- (sol);
\draw[arr] (cfg) -- (se);
\draw[arr] (cfg) -- (tu);
\draw[arr] (cfg) -- (sw);
\draw[arr] (cfg) -- (te);
\draw[arr] (rp)  -- (dyn);
\draw[arr] (rp)  -- (sol);
\draw[arr] (rp)  -- (se);
\draw[arr] (rp)  -- (tu);
\draw[arr] (rp)  -- (sw);
\draw[arr] (se)  -- (mc);
\draw[arr] (dyn) -- (mc);
\draw[arr] (sol) -- (mc);
\draw[arr] (gs)  -- (mc);
\draw[arr] (sw)  -- (mc);
\draw[arr] (te)  -- (mc);
\draw[arr] (tu)  -- (mc);
\draw[arr] (sp)  -- (mc);
\draw[arr] (mc)  -- (sim);
\draw[arr] (sim) -- (ft);
\draw[arr] (sim) -- (sc);
\end{tikzpicture}
\caption{Module dependency graph}
\end{figure}

\section{Per-Step Data Flow (Phase C)}

\begin{enumerate}
    \item \textbf{MuJoCo step} $\to$ \textbf{StateEstimator}: raw sensor data.
    \item \textbf{StateEstimator} $\to$ \textbf{RobotState}: structured state.
    \item \textbf{TerrainEstimator.update(state)}: IIR terrain height.
    \item \textbf{TrotGait.schedule(t, dt, K)}: $(K\times4)$ contact bool.
    \item \textbf{TrotGait.contact\_at(t)} $\to$ lift-off check:
          if transition detected, \textbf{SwingController.notify\_lift()}.
    \item \textbf{RigidBodyDynamics.Ac/Bc} for each $k$ $\to$ \textbf{ZOH} $\to$
          $(\mat{A}_d^{(k)}, \mat{B}_d^{(k)})$.
    \item \textbf{ConvexMPC.condense} $\to$ $(\mat{A}_\text{qp}, \mat{B}_\text{qp})$.
    \item \textbf{make\_reference} $\to$ $\mat{X}_\text{ref}$.
    \item \textbf{ConvexMPC.cost} $\to$ $(\mat{H}, \vect{g})$.
    \item \textbf{ConvexMPC.constraints} $\to$ $(\mat{C}, \vect{lb}, \vect{ub})$.
    \item \textbf{ConvexMPC.solve} $\to$ $\mat{U}_\text{opt}$.
    \item GRF smooth + re-seed $\to$ $\vect{f}_{1:4}$.
    \item \textbf{grf\_to\_joint\_torques}: $\vect{\tau}_\text{stance}$.
    \item \textbf{SwingController.compute\_torques}: $\vect{\tau}_\text{swing}$.
    \item Blend stance/swing torques by schedule.
    \item \textbf{leg\_to\_ctrl} (reorder) $\to$ clip $\to$ \texttt{data.ctrl}.
    \item \textbf{mj\_step()}.
\end{enumerate}

% ══════════════════════════════════════════════════════════════════════════════
\chapter{Bug Fixes and Design Decisions}
% ══════════════════════════════════════════════════════════════════════════════

\begin{longtable}{p{0.4cm}p{3.2cm}p{5cm}p{5.5cm}}
\caption{Complete bug fix and design decision log (19 items)}\\
\toprule
\textbf{\#} & \textbf{Location} & \textbf{Symptom / Root Cause} & \textbf{Fix} \\
\midrule
\endfirsthead
\multicolumn{4}{l}{\textit{(continued)}}\\
\toprule
\textbf{\#} & \textbf{Location} & \textbf{Symptom / Root Cause} & \textbf{Fix} \\
\midrule
\endhead
\bottomrule
\endfoot

1 & torque\_utils.py & Torque sign error: robot pushed itself into the ground.
    MPC gives GRF = upward force; joint must push \emph{down}.
  & Changed $+\mat{J}^\top\vect{f}$ to $-\mat{J}^\top\vect{f}$ \\
2 & state\_estimator.py & Phase A$\to$B boundary: efc\_force drops to 0 for $\sim$300\,ms.
    Contact reports NNNN; MPC sees all-swing, generates 3$\times$ overforce.
  & Height-based override: force contacts=True when $p_z > 0.20\,\text{m}$ \\
3 & config.py & Contact threshold $\tau_\text{cf} = 1\,\text{N}$ — below MuJoCo solver noise.
  & Raised to $10\,\text{N}$ \\
4 & mpc\_controller.py & Phase C reset pz\_des to NOMINAL\_HEIGHT (0.32\,m); body drops 18\,mm before first swing.
  & Keep pz\_des = settled height (0.338\,m) throughout \\
5 & gait\_scheduler.py & Gait phase measured from $t=0$; mid-cycle state at Phase C activation.
  & Anchor gait phase to activation time $t_0$ via \texttt{activate(t)} \\
6 & mpc\_controller.py & Q weights 10--50$\times$ too large (roll=100, pitch=200, pz=500); aggressive corrective forces destabilise trot.
  & Matched to CasADi reference values \\
7 & swing\_controller.py & Raibert overcorrection: $(t_\text{pred} + k_v)\,v_\text{actual}$ applied as one term; 3.5$\times$ too large.
  & Separated nominal drift ($v_\text{cmd}\,t_\text{pred}$) from error correction ($k_v\,(v_\text{actual}-v_\text{cmd})$) \\
8 & mpc\_controller.py & cmd\_vx/vy in body frame; Raibert needs world frame.
  & Rotate by yaw: $[v_x^w, v_y^w]^\top = \mat{R}_z(\psi)[v_x^\text{cmd}, v_y^\text{cmd}]^\top$ \\
9 & swing\_controller.py & Foot velocity used only 3-DOF leg Jacobian; ignored body rolling/pitching.
  & Use full $(3\times n_v)$ Jacobian: $\vect{v}_\text{foot} = \mat{J}_\text{full}\,\dot{\vect{q}}$ \\
10 & dynamics.py & Moment arm rotated into support-plane frame while forces stayed in world frame; physically wrong cross product.
   & Keep $\vect{r}_i$ in world frame; rotate only friction cones \\
11 & mpc\_solver.py & GRF smoothing was a no-op (result not stored back).
   & Fixed: $\texttt{self.\_grf} = \beta\,\vect{f}_\text{new} + (1-\beta)\,\texttt{self.\_grf}$ \\
12 & mpc\_controller.py & Terrain mode not reset in \texttt{simulation.reset()}; stair mode bled into viewer run.
   & Added \texttt{terrain.set\_mode(``flat'')} to \texttt{reset()} \\
13 & swing\_controller.py & Early touchdown fired on buffer bleed (20\,ms post-liftoff) in flat mode.
   & Guard 1: early touchdown disabled in flat mode \\
14 & swing\_controller.py & Early touchdown at 20\% swing progress froze foot on riser face, far short of tread.
   & Guard 2: require $>60\%$ swing progress before early touchdown \\
15 & terrain\_estimator.py & Smooth height pre-ramp generates phantom $+8^\circ$ pitch on flat approach.
   & Guard: return 0.0 when $x < 2.00\,\text{m}$ \\
16 & mpc\_controller.py & stairs\_Q not activated until $x > 2.00\,\text{m}$; first step hit with base Q (py=1, vy=2).
   & Activate stairs\_Q as soon as terrain mode switches \\
17 & mpc\_controller.py & pitch\_ref IIR $\alpha=0.5$ caused 3--6$^\circ$ lag; MPC generated anti-pitch torques opposing climb.
   & Changed to $\alpha=0.3$ (70\% new, 30\% old) \\
18 & mpc\_controller.py & Symmetric pz rate limiter $0.05\,\text{m/s}$: 30\,mm riser took 600\,ms; sustained overforce + roll crash.
   & Asymmetric: $0.12\,\text{m/s}$ up, $0.04\,\text{m/s}$ down \\
19 & mpc\_controller.py & GRF seed threshold 50\% insufficient; 100\,ms underforce window on new stance foot.
   & Raised to 70\% seed threshold \\
\end{longtable}

% ══════════════════════════════════════════════════════════════════════════════
\chapter{Numerical Summary}
% ══════════════════════════════════════════════════════════════════════════════

\begin{table}[H]
\centering
\caption{Key numerical parameters}
\begin{tabular}{ll}
\toprule
\textbf{Parameter} & \textbf{Value} \\
\midrule
Physics timestep       & $2\,\text{ms}$ (500\,Hz) \\
MPC solve rate         & $30\,\text{ms}$ (33\,Hz) \\
MPC horizon $K$        & 10 steps \\
Gait period $T$        & $0.40\,\text{s}$ (2.5\,Hz) \\
Duty cycle $d$         & 50\% \\
Stance time            & $0.20\,\text{s}$ \\
Swing time             & $0.20\,\text{s}$ \\
State dimension $n_s$  & 13 \\
Input dimension $n_u$  & 12 (4 feet $\times$ 3 forces) \\
QP size (stacked)      & $H \in \R^{120\times120}$ \\
Regularisation $\alpha$ & $3\times10^{-4}$ \\
Friction coefficient $\mu$ & 0.70 \\
Foot normal force range & $[5, 200]\,\text{N}$ \\
Torque limit per joint & $33.5\,\text{N\,m}$ \\
Phase A duration       & $1.5\,\text{s}$ \\
Phase B duration       & $2.0\,\text{s}$ \\
Phase C warmup         & $0.30\,\text{s}$ \\
Nominal CoM height     & $0.32\,\text{m}$ \\
Settled CoM height     & $\approx 0.338\,\text{m}$ \\
Collapse threshold     & $0.20\,\text{m}$ \\
Trot ready threshold   & $0.30\,\text{m}$ \\
Contact force threshold & $10\,\text{N}$ \\
GRF blend (flat)       & $\beta = 0.75$ \\
GRF blend (stairs)     & $\beta = 0.65$ \\
Stair approach x       & $1.80\,\text{m}$ \\
Flat Raibert $k_{v,x}$ & $0.4T = 0.16\,\text{m\,s}$ \\
Flat Raibert $k_{v,y}$ & $0.05T = 0.02\,\text{m\,s}$ \\
Swing clearance (flat) & $0.08\,\text{m}$ \\
Swing clearance (stairs) & $0.08\to0.18\,\text{m}$ (adaptive) \\
pz rate up             & $0.12\,\text{m/s}$ \\
pz rate down           & $0.04\,\text{m/s}$ \\
Pitch IIR $\alpha$     & $0.30$ (70\% new) \\
Tread safety margin    & $0.04\,\text{m}$ \\
Early touchdown guard  & $>60\%$ swing progress \\
\bottomrule
\end{tabular}
\end{table}

% ══════════════════════════════════════════════════════════════════════════════
\chapter{Conclusions}
% ══════════════════════════════════════════════════════════════════════════════

\section{Summary}

This project implements a complete convex MPC trot controller for the Unitree
Go2 quadruped robot in MuJoCo simulation.  The controller:

\begin{enumerate}
    \item Formulates the stance-force problem as a convex QP (13 states,
          12 inputs, $K=10$ horizon) solvable in real time.
    \item Uses a three-phase state machine (PD settle $\to$ all-stance MPC
          $\to$ trot MPC) to ensure smooth transitions.
    \item Implements a minimum-jerk swing trajectory with Raibert touchdown
          placement for forward progression.
    \item Extends to stair climbing through terrain-aware references
          (adaptive $p_z$, pitch reference), per-foot stair snapping,
          adaptive swing clearance, support-plane friction cones, and
          early touchdown detection.
    \item Documents 19 bug fixes across 4+ debugging iterations, covering
          signal processing (contact detection), control theory (Raibert
          formulation, frame conventions), and numerical stability (GRF
          smoothing, rate limiters).
\end{enumerate}

\section{Limitations and Future Work}

\begin{itemize}[noitemsep]
    \item The SRB model ignores leg mass ($\approx 10\%$ of total), which
          limits accuracy during dynamic swing.  A whole-body MPC would improve
          this at significant computational cost.
    \item The trot gait uses a fixed period.  Adaptive gait timing (speed-dependent
          $T$) would improve efficiency over a wider velocity range.
    \item The terrain estimator uses a hardcoded stair geometry.  A general
          elevation map (e.g.\ from depth sensors) would enable uneven terrain.
    \item The friction coefficient $\mu = 0.7$ is fixed.  Online friction
          estimation from force/slip measurements would improve robustness.
    \item The controller has been validated only in simulation.  Sim-to-real
          transfer requires addressing actuator delays, sensor noise, and
          unmodelled flexibility.
\end{itemize}

% ══════════════════════════════════════════════════════════════════════════════
\bibliographystyle{plain}
\begin{thebibliography}{9}

\bibitem{diCarlo2018}
J.~Di Carlo, P.~M.~Wensing, B.~Katz, G.~Bledt, and S.~Kim,
``Dynamic locomotion in the MIT Cheetah 3 through convex model-predictive control,''
in \textit{Proc. IEEE/RSJ Int. Conf. Intelligent Robots and Systems (IROS)},
pp. 7440--7447, 2018.

\bibitem{raibert1986}
M.~H.~Raibert,
\textit{Legged Robots That Balance}.
MIT Press, Cambridge, MA, 1986.

\bibitem{unitreeGo2}
Unitree Robotics,
``Go2 Robot Dog,''
\texttt{https://www.unitree.com/go2/}, 2023.

\bibitem{mujoco}
E.~Todorov, T.~Erez, and Y.~Tassa,
``MuJoCo: A physics engine for model-based control,''
in \textit{Proc. IEEE/RSJ IROS}, pp. 5026--5033, 2012.

\end{thebibliography}

\end{document}
